\section{Initial data and sampling}\label{app:sampling}

Coupled Einstein-Vlasov evolutions require constraint-satisfying initial slices alongside discrete particle ensembles that sample the continuum phase-space distribution. Each self-gravitating run starts on a conformally flat, time-symmetric slice where superposed flat-space Poisson potentials $u_c$ for conformal sources $\tilde{\rho}_c$ solve $\Ham=0$ via Eq.~(\ref{eq:initial_superposition}). Around each component center $\bm{x}_c$ at isotropic distance $\rbar_c=|\bm{x}-\bm{x}_c|$, four component families enter Eq.~(\ref{eq:initial_superposition}). Tabulated spherical profiles interpolate $u_c(\rbar_c)=\psi_c(\rbar_c)-1$ on a uniform grid with four-point Lagrange polynomials and follow Eq.~(\ref{eq:puncture_potential}) outside with the component ADM mass. Vacuum black hole punctures evaluate Eq.~(\ref{eq:puncture_potential}) directly, while spherical Gaussian components of conformal mass $m_c$ and width $\sigma_c$ obey
\begin{align}\label{eq:app_gaussian}
\tilde{\rho}_c(\rbar_c)&=\frac{m_c}{(2\pi\sigma_c^2)^{3/2}}\exp\!\left(-\frac{\rbar_c^2}{2\sigma_c^2}\right),\nonumber\\
u_c(\rbar_c)&=\frac{m_c}{2\rbar_c}\,\mathrm{erf}\!\left(\frac{\rbar_c}{\sqrt{2}\,\sigma_c}\right).
\end{align}
A homogeneous prolate spheroid of ADM mass $M$, polar semi-axis $b$ along $z$, eccentricity $e$, equatorial semi-axis $a_{\mathrm{eq}}=b\sqrt{1-e^2}$, and focal half-distance $d_{\mathrm{f}}=eb$ has closed-form fields \cite{Shapiro:1991zza,Yoo:2016kzu}
\begin{align}\label{eq:app_spheroid}
\tilde{\rho}_c(\bm{x})&=\frac{3M}{4\pi a_{\mathrm{eq}}^2b},\qquad\frac{x^2+y^2}{a_{\mathrm{eq}}^2}+\frac{z^2}{b^2}\le1,\nonumber\\
u_c(\bm{x})&=\frac{3M}{8}\Bigl[I_0(\lambda)-(x^2+y^2)J_{\mathrm{eq}}(\lambda)-z^2J_z(\lambda)\Bigr],\nonumber\\
I_0(\lambda)&=\frac{2}{d_{\mathrm{f}}}\,\mathrm{arctanh}\!\left(\frac{d_{\mathrm{f}}}{b_\lambda}\right),\nonumber\\
J_{\mathrm{eq}}(\lambda)&=\frac{b_\lambda}{d_{\mathrm{f}}^2(a_{\mathrm{eq}}^2+\lambda)}-\frac{I_0(\lambda)}{2d_{\mathrm{f}}^2},\nonumber\\
J_z(\lambda)&=\frac{I_0(\lambda)}{d_{\mathrm{f}}^2}-\frac{2}{d_{\mathrm{f}}^2b_\lambda},\nonumber\\
1&=\frac{x^2+y^2}{a_{\mathrm{eq}}^2+\lambda}+\frac{z^2}{b^2+\lambda}\qquad(\lambda>0).
\end{align}
The auxiliary length $b_\lambda=\sqrt{b^2+\lambda}$ sets $\lambda=0$ inside the spheroid, while outside the boundary $\tilde{\rho}_c$ vanishes and $\lambda>0$ solves the last relation. Every initial slice sets vanishing shift $\beta^i=0$, with initial lapse $\alpha=1$ for cold collapse and $\alpha=\psi^{-2}$ for vacuum punctures. Isolated clusters use their static equilibrium lapse joined outside to $(1-u_c)(1+u_c)^{-1}$, whereas the composite cluster uses $\alpha=2\psi^{-1}-1$. Since $\sqrt{\gamma}=\psi^6$, $N$ equal-mass macroparticles sampling proper rest-mass density $\rho_0$ and continuum rest mass $M_0$ have coordinate number density
\begin{align}\label{eq:app_coord_density}
\mathcal{N}(\bm{x})&=\frac{N}{M_0}\,\rho_0(\bm{x})\psi(\bm{x})^6,\nonumber\\
m_p&=\frac{M_0}{N},\qquad M_0=\int\rho_0(\bm{x})\psi(\bm{x})^6\,\dd^3x.
\end{align}
For the cold prolate spheroid ($u_{i,p}=0$, $\rho_0\psi^6=\tilde{\rho}_c\psi$), one quarter of the $N$ particles are drawn inside the boundary by rejection against $\psi(\bm{x})\psi(\bm{0})^{-1}$ and completed by $(x,y,-z)$ and $(-x,-y,\pm z)$. Gaussian components draw positions around $\bm{x}_c$ by randomized stratification in radius and angles. In multi-component superpositions, each component multiplies its isolated particle masses by $\psi\psi_c^{-1}$ ($\psi_c=1$ and rest mass $m_cW_c^{-1}$ at bulk Lorentz factor $W_c$) and lowers orthonormal momenta with $\psi^2$.

The homogeneous dust ball of mass $M$, areal radius $R_0$, and density $\rho_0$ (Sec.~\ref{sec:dust}) starts from rest ($u_{i,p}=0$) with isotropic radius $\rbar_0$ and interior conformal factor \cite{Staley:2011ss}
\begin{align}\label{eq:app_dust_psi}
\rbar_0&=\frac{R_0}{2}\!\left(1-\frac{M}{R_0}+\sqrt{1-\frac{2M}{R_0}}\right),\nonumber\\
\psi(\rbar)&=\sqrt{\frac{\bigl(1+\sqrt{1-2M/R_0}\bigr)\rbar_0R_0^2}{2\rbar_0^3+M\rbar^2}}\qquad(\rbar\le\rbar_0),
\end{align}
matched to Eq.~(\ref{eq:puncture_potential}) outside. The cube $[-\rbar_0,\rbar_0]^3$ is partitioned into $n_{\mathrm{lat}}^3$ cubic cells of spacing $h=2\rbar_0n_{\mathrm{lat}}^{-1}$ aligned with the initial matter-level mesh, with $n_{\mathrm{lat}}$ in the parameter table. Macroparticles occupy cell centers with $\rbar_p=|\bm{x}_p|\le\rbar_0$ for noise-free interior deposition, and the published rule \cite{Wu:2026nrpic} assigns proper cell masses
\begin{equation}\label{eq:app_dust_published}
m_p^{\mathrm{pub}}=\rho_0\psi(\rbar_p)^6h^3.
\end{equation}
The count of lattice sites inside a sphere differs from its volume in cells $h^3$ by a boundary fraction of order $n_{\mathrm{lat}}^{-1}$, so summed $m_p^{\mathrm{pub}}$ exceeds $M_0$ on small lattices. Our surface correction leaves interior masses unchanged for $\rbar_p\le\rbar_0-h$ and rescales the outermost layer $\rbar_0-h<\rbar_p\le\rbar_0$ of width $h$ by one common factor so all particle masses sum to $M_0$,
\begin{equation}\label{eq:app_dust_surface}
m_p=m_p^{\mathrm{pub}}\,\frac{M_0-\sum_{\rbar_{p'}\le\rbar_0-h}m_{p'}^{\mathrm{pub}}}{\sum_{\rbar_{p'}>\rbar_0-h}m_{p'}^{\mathrm{pub}}}.
\end{equation}

Monoenergetic equilibrium clusters \cite{1985ApJ...298...34S} populate a single conserved energy per unit rest mass $E_0=\alpha_{\mathrm{s}}$ in Eq.~(\ref{eq:cluster_sequence}) with isotropic momentum directions. In dimensionless areal radius $\tilde{r}$, squared lapse ratio $y(\tilde{r})$ with $y(0)=y_c$, and enclosed gravitational and rest-mass variables $\tilde{m}$ and $\tilde{m}_0$, the structure equations and isotropic-radius map take the form
\begin{align}\label{eq:app_structure}
\zeta&=\frac{1-y_c}{3},\qquad g_y=\frac{y^{-1}-1}{y_c^{-1}-1},\nonumber\\
\frac{\dd y}{\dd\tilde{r}}&=\frac{2\zeta y}{\tilde{r}^2}\,\frac{\tilde{m}+\zeta\tilde{r}^3(y/y_c)^{-1/2}g_y^{3/2}}{1-2\zeta\tilde{m}/\tilde{r}},\nonumber\\
\frac{\dd\tilde{m}}{\dd\tilde{r}}&=\tilde{r}^2\!\left(\frac{y}{y_c}\right)^{\!-3/2}g_y^{1/2},\nonumber\\
\frac{\dd\tilde{m}_0}{\dd\tilde{r}}&=\frac{\tilde{r}^2y_c^{1/2}(y/y_c)^{-1}g_y^{1/2}}{\sqrt{1-2\zeta\tilde{m}/\tilde{r}}},\nonumber\\
\frac{\dd\ln\rbar}{\dd\ln\tilde{r}}&=\frac{1}{\sqrt{1-2\zeta\tilde{m}/\tilde{r}}}.
\end{align}
Near the cluster surface where density vanishes, the variable $\sqrt{1-y}$ keeps every derivative analytic up to $y(\tilde{r}_{\mathrm{s}})=1$. Outward integration from $\tilde{r}=0$ ($\tilde{m}=\tilde{m}_0=0$) yields surface values $(\tilde{r}_{\mathrm{s}},\tilde{m}_{\mathrm{s}},\tilde{m}_{0,\mathrm{s}})$ and interior profiles matched at $\rbar_{\mathrm{s}}$ from Eq.~(\ref{eq:app_dust_psi}),
\begin{align}\label{eq:app_cluster_profiles}
R&=\frac{M\tilde{r}_{\mathrm{s}}}{\zeta\tilde{m}_{\mathrm{s}}},\qquad M_0=\frac{M\tilde{m}_{0,\mathrm{s}}}{\tilde{m}_{\mathrm{s}}},\nonumber\\
E_0&=\sqrt{1-\frac{2M}{R}},\qquad\alpha_{\mathrm{eq}}(\rbar)=E_0\sqrt{y},\nonumber\\
\psi_{\mathrm{eq}}(\rbar)&=\sqrt{\frac{R\tilde{r}}{\tilde{r}_{\mathrm{s}}\rbar}},
\end{align}
with enclosed masses $m_{\mathrm{g}}=M\tilde{m}\tilde{m}_{\mathrm{s}}^{-1}$, $m_0=M_0\tilde{m}_0\tilde{m}_{0,\mathrm{s}}^{-1}$, and $4\pi\rbar^2\rho_0\psi_{\mathrm{eq}}^6\,\dd\rbar=\dd m_0$. For $N$ equal-mass macroparticles ($p=1,\dots,N$) and five coordinates $d\in\{1,\dots,5\}$, randomized stratification draws
\begin{equation}\label{eq:app_stratified}
U_{pd}=\frac{\pi_d(p)+\xi_{pd}}{N},
\end{equation}
where independent random permutations $\pi_d$ map $\{1,\dots,N\}$ onto $\{0,\dots,N-1\}$ and $\xi_{pd}$ is uniform on $[0,1)$. The cumulative rest-mass fraction $F(\rbar)=m_0(\rbar)M_0^{-1}$ sets each radial coordinate through $\rbar_p=F^{-1}(U_{p1})$. Position angles $(\cos\theta_p,\phi_p)=(2U_{p2}-1,2\pi U_{p3})$ define the radial unit vector $\hat{\bm{r}}_p$ and $\bm{x}_p=\rbar_p\hat{\bm{r}}_p$, while momentum angles $(\cos\theta_{u,p},\phi_{u,p})=(2U_{p4}-1,2\pi U_{p5})$ define an isotropic unit vector $\hat{\bm{u}}_p$. At radius $\rbar_p$, the monoenergetic condition $E_0=\alpha_{\mathrm{eq}}(\rbar_p)W_p$ and Eq.~(\ref{eq:lorentz_factor}) fix the local orthonormal momentum magnitude $u_{\mathrm{eq}}(\rbar_p)$ and covariant momentum,
\begin{align}\label{eq:app_cluster_momentum}
u_{\mathrm{eq}}(\rbar)&=\sqrt{\frac{E_0^2}{\alpha_{\mathrm{eq}}(\rbar)^2}-1},\nonumber\\
u_{i,p}&=\psi_{\mathrm{eq}}(\rbar_p)^2u_{\mathrm{eq}}(\rbar_p)\hat{u}_{i,p}.
\end{align}
Plain Monte Carlo sampling, used as the alternative in the particle-number study, replaces $U_{pd}$ in Eq.~(\ref{eq:app_stratified}) with independent uniform draws on $[0,1)$. Because the mass-weighted mean momentum of a finite sample does not vanish and makes the cluster drift, that mean is subtracted from every particle,
\begin{equation}\label{eq:app_zero_momentum}
u_{i,p}\leftarrow u_{i,p}-\frac{\sum_{p'=1}^{N}m_{p'}u_{i,p'}}{\sum_{p'=1}^{N}m_{p'}}.
\end{equation}

Along the equilibrium sequence, an inward velocity kick of fractional amplitude $\delta_v$ scales the radial momentum component as in Ref.~\cite{Wu:2026nrpic}, shifting $\hat{u}_{i,p}$ in Eq.~(\ref{eq:app_cluster_momentum}) to $\hat{u}_{i,p}-\delta_v\hat{r}_{i,p}$ before Eq.~(\ref{eq:app_zero_momentum}). For the homologous density perturbation $\delta_\rho$ of the near-critical model, equilibrium radii $\rbar_{\mathrm{eq},p}$ scale by a common factor $\Lambda_{\mathrm{h}}$ to $\rbar_p=\Lambda_{\mathrm{h}}\rbar_{\mathrm{eq},p}$ with $\bm{x}_p=\rbar_p\hat{\bm{r}}_p$. Unlike a perturbation applied on the unperturbed metric $\psi_{\mathrm{eq}}$, here $\psi(\rbar)$ is recomputed from the perturbed matter by a radial fixed-point solve of the Hamiltonian constraint,
\begin{align}\label{eq:app_perturbed_hc}
\Lambda_{\mathrm{h}}&=(1+\delta_\rho)^{-1/3},\nonumber\\
\frac{1}{\rbar^2}\frac{\dd}{\dd\rbar}\!\left(\rbar^2\frac{\dd\psi}{\dd\rbar}\right)&=-\frac{2\pi}{\Lambda_{\mathrm{h}}^3\psi(\rbar)}\,\rho_0\!\left(\frac{\rbar}{\Lambda_{\mathrm{h}}}\right)\psi_{\mathrm{eq}}\!\left(\frac{\rbar}{\Lambda_{\mathrm{h}}}\right)^{\!6}\bigl\langle W(\rbar)\bigr\rangle.
\end{align}
Outside $\Lambda_{\mathrm{h}}\rbar_{\mathrm{s}}$, $\psi(\rbar)$ joins $1+M'(2\rbar)^{-1}$ for perturbed ADM mass $M'$, while the interior lapse rescales $\alpha_{\mathrm{eq}}(\rbar\Lambda_{\mathrm{h}}^{-1})$ continuously to the exterior isotropic Schwarzschild lapse. Particle momenta at $\rbar_p$ are set from $E_0$ via $u_{\mathrm{eq}}(\rbar_{\mathrm{eq},p})$ lowered with $\psi_{\mathrm{eq}}(\rbar_p)^2$ before Eq.~(\ref{eq:app_zero_momentum}), while $\langle W(\rbar)\rangle$ averages Eq.~(\ref{eq:lorentz_factor}) over directions in the new metric $\psi(\rbar)^4\delta_{ij}$. An inward velocity kick with vanishing extrinsic curvature $K_{ij}=0$ leaves a first-order violation of the momentum constraint $\Mom_i=0$, which is not corrected.

The relativistic cluster with two collapsing overdensities superposes a monoenergetic envelope and two half-turn-symmetric cold Gaussian components on an equatorial circular orbit of radius $d_{\mathrm{orb}}$. At orbital radius $d_{\mathrm{orb}}$ enclosing envelope mass $m_{\mathrm{g}}(d_{\mathrm{orb}})$, leading-order circular motion gives
\begin{equation}\label{eq:app_circular_speed}
v_{\mathrm{orb}}^2=\frac{m_{\mathrm{g}}(d_{\mathrm{orb}})}{d_{\mathrm{orb}}}+\frac{m_c}{4d_{\mathrm{orb}}},
\end{equation}
which matches the speed quoted in Ref.~\cite{Wu:2026nrpic} only for a host cluster less than half as large as the stated radius. The cluster size in the published figure agrees with that smaller value, and the rebuild adopts this self-consistent host radius in the parameter table.
